
\subsection{\texorpdfstring{\textbf{
                  Структура модели Хопфилда}}
      {Структура модели Хопфилда}}

В начале возникают вопросы. Как выглядит нейронная сеть Хопфилда?
Как обучается? Как используется?

Классическая модель сети Хопфилда состоит из состояния нейронов -- x -
вектор длины N и весов - матрицы связей от нейрона i к нейрону j - T --
размера N на N. Важно, что матрица весов T непрерывна. Учится на
паттернах - векторах длины N \(p^{k}\).

Нейроны могут хранить как двоичную информацию, так и вещественные числа,
в зависимости от выбора динамики и правила обучения.

Теоретическая основа следующая: обучение модели смещает функционал
энергии так, чтобы при входном паттерне динамика минимизировала его к
одной из энергий при паттерне, на котором она была обучена.

Очень хорошую аналогию можно провести с гравитацией. Обучение --
размещение планет в системе, динамика -- гравитация, применение --
размещённая планета притягивается к планете, которая больше всего её
притягивает.

Из аналогии также вытекает проблема, из-за которой сети Хопфилда плохо
работают с коррелированными паттернами. Если паттерны похожи друг на
друга, то нейросеть рискует выдать не тот паттерн или склонна выдавать
паттерн, которого изначально не было при обучении. Второй случай
возникает, когда совместная гравитация двух и более планет создаёт один
минимум, а не несколько.

Это также заметно и в работе памяти человека. Образы чаще всего размыты
и часто смешиваются. Возможно, эту проблему нельзя решить точно, не
использовав больше ресурсов. Известно, что эволюционно выгоднее
затрачивать меньше ресурсов. Проблема фундаментальна - серое пятно нельзя
идентифицировать точно. То есть мы можем различать только паттерны, которые
не сильно отличны от исходного.
